\documentclass[10pt,a4paper]{amsart}
\usepackage[margin=19mm]{geometry}
\usepackage{amsmath,amssymb,mathtools}
\usepackage[scr=rsfs,cal=euler]{mathalfa}
\usepackage{xfrac}
\usepackage[unicode,hidelinks,bookmarks=false]{hyperref}
\newcommand{\ie}{i.e.\ }
\newcommand{\cf}{cf.\ }
\newcommand{\resp}{respectively\ }
\newcommand{\Subsection}{\subsection}
\providecommand{\nobreakdash}{\nobreak\hbox{-}}
\setlength{\emergencystretch}{2em}
\multlinegap0pt
\usepackage[shortlabels]{enumitem}
\usepackage{amssymb}
\usepackage{mathtools}
\usepackage{xcolor}
\usepackage{tikz}
\newtheorem{corollary}{Corollary}
\newtheorem{proposition}{Proposition}
\usepackage{cleveref}
\crefname{corollary}{Corollary}{Corollarys}
\crefname{proposition}{Proposition}{Propositions}
\newcommand{\Ad}{\operatorname{Ad}}
\newcommand{\Lie}{\operatorname{Lie}}
\newcommand{\RR}{\mathbb R}
\newcommand{\SO}{\mathrm{SO}}
\newcommand{\SU}{\mathrm{SU}}
\title{\textbf{Quaternionic Reflections and Lie Generation in $\mathfrak g_2$}}
\author{Santiago Pineda Montoya \and Johan H. R\'{u}a Mu\~{n}oz \and Borut Jur\v{c}i\v{c} Zlobec}
\date{Source: arXiv:2609.15497v1 (CC BY 4.0)}
\begin{document}
\maketitle

\noindent\emph{Source and licence.} \textbf{Quaternionic Reflections and Lie Generation in $\mathfrak g_2$}. Version arXiv:2609.15497v1 (CC BY 4.0), distributed by arXiv under the Creative Commons Attribution 4.0 International licence, http://creativecommons.org/licenses/by/4.0/, as recorded in the arXiv licence metadata for this exact version and shown on the abstract page https://arxiv.org/abs/2609.15497v1. Full licence notice: \texttt{licenses/arxiv-2609.15497-CC-BY-4.0.txt}.\par\smallskip

\noindent\textit{Editorial scope.} Counted unit: the article's proposition prop:principal-reflection-fiber (source0.tex lines 1148--1156). The article's complete original proof of it is delivered with it (source0.tex lines 1157--1198, one \texttt{proof} environment of 42 lines). No mathematics is written, repaired or re-derived: the statement and the proof are the article's own lines, in the article's own notation, and no retained line is substituted or re-flowed.  Retained uncounted as the source-local context the proof needs: lines 389--389; lines 1084--1097.  Nothing was omitted from any retained span, so the package carries no omission record.  No retained line contains a citation, so the wrapper carries no bibliography and no bibliography entry.  No retained line contains a figure environment, an image-inclusion command or drawing code, so no image is delivered and the package declares no asset dependency.  Editorial formatting only, in addition: an AMS \texttt{amsart} preamble that restates the article's own declarations and macros used by the retained lines, namely the environments \texttt{corollary}, \texttt{proposition} on separate counters, together with the standard packages those lines need.  The retained environments print in this document as Corollary 1, Proposition 1 in order of appearance, whereas the article prints the same items with the numbering of its own full document.  The complete native source of the article as distributed in the arXiv e-print for this exact version is bundled unchanged as \texttt{sources/210421/source0.tex}.
\subsection{Compact maximal subalgebras}\label{subsec:compact-maximals}

\begin{corollary}[Spectral realizability of the principal obstruction]\label{cor:universal-spectral}
For a nonzero $Z\in\mathfrak g_2$, the following conditions are equivalent:
\begin{enumerate}[label=\textnormal{(\roman*)},leftmargin=9mm]
\item the positive rotation frequencies of $Z$ have ratio $1:2:3$;
\item $Z$ belongs to a principal $\mathfrak{su}(2)$ subalgebra of $\mathfrak g_2$;
\item there is $T\in\mathscr C$ such that $\Lie\{Z,TZT\}$ is a principal $\mathfrak{su}(2)$.
\end{enumerate}
For $Z=Z(q_a,q_b)\in\mathfrak k_A$ with $\alpha>0$, these conditions are equivalent to
\[
\frac{\beta}{\alpha}\in\left\{2,5,\frac13\right\}.
\]
Every reflection realizing (iii) for such a local $Z$ moves $H_A$.
Outside the principal spectrum, including $Z=0$, generation for any $T\in\mathscr C$ is equivalent to $\mathcal R_{30}(T,Z)>0$.
\end{corollary}

\begin{proposition}[The realizing-reflection fiber]\label{prop:principal-reflection-fiber}
The centralizer $C_{\rm p}$ is a two-dimensional torus.  For the fixed choice $Z=c\Ad_gK_{\rm p}$,
\[
\mathscr P_Z
 =\{g uT_\theta u^{-1}g^{-1}:u\in C_{\rm p},\quad 0<\theta<\pi/2\}.
\]
Every reflection on the right has a unique pair $(u,\theta)$.  This parametrization is a real-analytic diffeomorphism onto $\mathscr P_Z$ with the subspace topology inherited from $\mathscr C$.  The fiber is therefore an embedded real-analytic submanifold of dimension three and codimension five in the eight-dimensional reflection manifold $\mathscr C$.
If $Z\in\mathfrak k_A$, all these reflections move $H_A$.
\end{proposition}

\begin{proof}
Scaling and conjugation reduce the proof to $Z=K_{\rm p}$.
First compute its group centralizer.  The kernel line is $\RR e_1$.  An element centralizing $K_{\rm p}$ must preserve this line.
If it sent $e_1$ to $-e_1$, its restriction to $e_1^\perp$ would be conjugate-linear for $C_{e_1}$, since it preserves the cross product.
Commutation with $K_{\rm p}$ would then identify the signed complex spectrum $(2,1,-3)$ with its negative, which is impossible.
It therefore fixes $e_1$ and lies in its $\SU(3)$ stabilizer.
There $K_{\rm p}$ has three distinct complex eigenvalues, so its centralizer is the diagonal determinant-one unitary torus.  Hence $C_{\rm p}\cong(S^1)^2$.

The group centralizer of $\mathfrak s_{\rm p}$ in $G_2$ is trivial.  Indeed, its real seven-dimensional action is irreducible.
A commuting symmetric operator is scalar, and a commuting skew-symmetric operator in odd dimension has a nonzero kernel and must vanish by irreducibility.
Thus the real commutant consists of scalars; its orthogonal elements are $\pm I_7$, of which only $I_7$ lies in $\SO(7)$.
Every automorphism of $\mathfrak{su}(2)$ is inner and is realized by $S_{\rm p}$, so
\[
N_{G_2}(\mathfrak s_{\rm p})=S_{\rm p}.
\]
Every principal algebra containing $K_{\rm p}$ is a $C_{\rm p}$-conjugate of $\mathfrak s_{\rm p}$.
To see this, choose $h$ carrying $\mathfrak s_{\rm p}$ to such an algebra, using the compact conjugacy in \Cref{subsec:compact-maximals}.
Then $\Ad_{h^{-1}}K_{\rm p}\in\mathfrak s_{\rm p}$ has the same trace norm as $K_{\rm p}$.
Rotations in $S_{\rm p}$ are transitive on that sphere, so $\Ad_vK_{\rm p}=\Ad_{h^{-1}}K_{\rm p}$ for some $v\in S_{\rm p}$.
The element $hv$ centralizes $K_{\rm p}$ and carries $\mathfrak s_{\rm p}$ to the prescribed algebra.

If $T\in\mathscr P_{K_{\rm p}}$, it normalizes its generated principal algebra: conjugation by $T$ interchanges the two generators.
Choose $u\in C_{\rm p}$ with $\Lie\{K_{\rm p},TK_{\rm p}T\}=\Ad_u\mathfrak s_{\rm p}$, as in the preceding paragraph.
Then $u^{-1}Tu\in N_{G_2}(\mathfrak s_{\rm p})=S_{\rm p}$ is a nonidentity involution.
It is therefore a half-turn about an unoriented axis in $S_{\rm p}\cong\SO(3)$.  Write its unit axis in the normalization $(K_{\rm p},P_{\rm p},Q_{\rm p})$ as
$\xi_1K_{\rm p}+\xi_2P_{\rm p}+\xi_3Q_{\rm p}$, where $\sum\xi_j^2=1$.
Its conjugate of $K_{\rm p}$ is independent of $K_{\rm p}$ exactly when $0<|\xi_1|<1$.
Choose the sign of the axis so that $\xi_1>0$ and rotate its transverse direction by $\exp(tK_{\rm p})\in C_{\rm p}$.
The axis becomes $\cos\theta K_{\rm p}+\sin\theta P_{\rm p}$ with $0<\theta<\pi/2$.
Conversely, every such half-turn generates $\mathfrak s_{\rm p}$ with $K_{\rm p}$ and is a quaternionic reflection, by the spin-three weights.
This proves the asserted equality of sets.

The angle is recovered from
\[
\frac{\langle K_{\rm p},TK_{\rm p}T\rangle_{\rm tr}}{\|K_{\rm p}\|_{\rm tr}^2}=\cos(2\theta).
\]
If two centralizer elements give the same reflection at this angle, their quotient commutes with $K_{\rm p}$ and $T_\theta$.
It therefore commutes with $K_{\rm p}$ and $T_\theta K_{\rm p}T_\theta$, hence with all of $\mathfrak s_{\rm p}$, and is the identity.
The parametrization is thus injective.  Its torus-orbit differential is injective by the same stabilizer argument; the angle has nonzero derivative in the displayed scalar on the open interval.
It is an analytic immersion.  On every compact subinterval the domain is compact, and the displayed scalar recovers $\theta$ continuously in the subspace topology of $\mathscr C$.  It follows that the map is an embedding with analytic inverse onto its image.  Its codimension is $\dim\mathscr C-3=8-3=5$.
The last assertion follows from \Cref{cor:universal-spectral}.
\end{proof}

\end{document}
